% Requires booktabs. Detailed protocol: stateful_appendix.tex.
\subsection{Controlled training with stateful dependency tools}
\label{sec:stateful-training}

We test whether the label estimators support actual policy updates in a
constructed release-repair environment. A component build records the current
artifact of its parent; changing an upstream component can therefore leave
downstream builds, attestations, and published artifacts stale. The cheap
checker inspects only published local versions, whereas the reference checker
executes all dependency and provenance constraints. These are deterministic
tools and a NumPy softmax policy, not a language model or a public software
benchmark.

We run all six methods for 300 on-policy updates, with groups of eight
trajectories, horizons $H\in\{12,24,48\}$, and ten independent training seeds
per setting. All methods infer reference failure for free when the necessary
cheap check fails. If $M$ candidate labels remain, each partial method purchases
$M/2$ labels in expectation. The completion probabilities and sequential
survival schedule are fixed before training. Evaluation uses 256 independent
reporting episodes per checkpoint, paired across methods within a seed.

\begin{table}[t]
\centering
\small
\setlength{\tabcolsep}{4pt}
\begin{tabular}{lrrr}
\toprule
Method & $H=12$ & $H=24$ & $H=48$ \\
\midrule
Full-label update & $94.06\pm0.61$ & $83.44\pm1.03$ & $57.97\pm0.96$ \\
Cheap labels only & $17.42\pm3.58$ & $19.84\pm2.92$ & $3.59\pm0.95$ \\
Whole-group residual & $93.48\pm0.71$ & $83.24\pm1.05$ & $55.27\pm1.16$ \\
Sequential completion & $93.55\pm0.63$ & $83.05\pm0.88$ & $56.60\pm1.72$ \\
Reward HT then normalize & $86.84\pm0.47$ & $68.91\pm1.07$ & $32.23\pm1.24$ \\
Partial replacement & $91.13\pm0.58$ & $75.08\pm0.68$ & $29.96\pm1.76$ \\
\bottomrule
\end{tabular}
\caption{Reference-check success (\%, mean $\pm$ standard error across ten
training seeds) after the same 2,400 training trajectories per run.
Purchased-label counts differ. The initial mean success rates were
$40.08\%$, $18.55\%$, and $4.96\%$, respectively.}
\label{tab:stateful-training}
\end{table}

Sequential completion (SC) and the strong whole-group residual baseline reach similar
mean endpoints (Table~\ref{tab:stateful-training}); all three paired,
pointwise 95\% intervals for their difference include zero. Thus these runs do
not establish a distinct benefit from the fixed sequential schedule. Nor does
an interval containing zero establish equivalence to full-label training.
Cheap-label optimization can instead improve the incomplete checker while
retaining poor reference success: at $H=48$, its endpoint cheap success is
$92.38\%$, versus $3.59\%$ under the reference checker.

Across the 30 runs for each method, SC purchases 30,908 training labels,
compared with 62,508 for full-label updates and 30,832 for whole-group residual
correction. This query reduction is not a measured speedup: training takes
29.08, 19.79, and 21.35 seconds, respectively, with these inexpensive checks.
The scripted initialization already executes the known workflow greedily;
this experiment measures stochastic execution reliability, not discovery of
an unknown repair procedure. Further protocol, accounting, and inferential
limits appear in Appendix~\ref{app:stateful-training}.
